\documentclass{dreamClass}

\title{\heiti\textbf{某门课 \quad 某次作业}}
\author{123456789 赵志辉\thanks{\href{mailto:test@example.com}{test@example.com}}}
\affil{九乡河文理学院}
\date{\today}

\begin{document}
\maketitle
\thispagestyle{empty}

\section{喵}
下面是带圆圈的编号列表：
\begin{enumerate}[label=\large\protect\textcircled{\small\arabic*}]
    \item 喵
    \item 喵
\end{enumerate}

\autoref{pic:desktop}是一幅插图。
\begin{figure}
    \centering
    \includegraphics[width = 0.5\textwidth]{screenshot.jpg}
    \caption{一张截图}\label{pic:desktop}
\end{figure}

\autoref{code:junk}是作为浮动体的代码块：
\begin{listing}
    \caption{凑字数用的代码\label{code:junk}}
    \begin{minted}{cpp}
    int main() {}
    \end{minted}
\end{listing}

\autoref{code:junk-again}是并非浮动体的代码块：
\begin{codeblock}
    \caption{复读机\label{code:junk-again}}
    \begin{minted}{cpp}
    int main() {}
    \end{minted}
\end{codeblock}

这是行内的代码片段：\mintinline{cpp}{wstring price = L"九磅十五便士"}。

\autoref{code:engine}是使用文件作为代码块的内容：
\begin{codeblock}
    \caption{\texttt{Engine.hpp}的内容\label{code:engine}}
    \inputminted{cpp}{../code/engine/Engine.hpp}
\end{codeblock}

下面是一些数学式：
\[
    \begin{align*}
        P(x) & = \int_{0}^{x} \beta e^{ - \beta x} \dif x \\
             & = \int_{0}^{\beta x} e^{ - t} \dif t       \\
             & = - e^{ - u} \Big|_{0}^{\beta x}           \\
             & = 1 - e^{ - \beta x}\text{。}
    \end{align*}
\]
\[
    \begin{aligned}
        \begin{array}{cl}
            \text{maximize} & H = \bm{p}^T\bm{q} \\
            \text{ s.t. }   & \bm{p}^T\bm{1} - 1 = 0\text{。}
        \end{array}
    \end{aligned}
\]
考虑\(\mathbf{H}_{i j}=\frac{\partial^{2} f}{\partial x_{i} \partial x_{j}}\)。
又\(\operatorname{Pr}(X \geq 10^3 + 2000 | X \geq 2000) = \operatorname{Pr}(X \geq 10^3)\)，

Matlab的在线手册告诉我们\footnote{见\href{https://www.mathworks.com/help/matlab/ref/power.html}{https://www.mathworks.com/help/matlab/ref/power.html}}：
\end{document}
